\begin{thebibliography}{99}
\footnotesize\itemsep=0pt

\bibitem{A}
Andrews G.E., On a transformation of bilateral series with applications,
  \href{https://doi.org/10.2307/2036642}{\textit{Proc. Amer. Math. Soc.}}
  \textbf{25} (1970), 554--558.

\bibitem{AB}
Andrews G.E., Berndt B.C., Ramanujan's lost notebook. {P}art~{II},
  \href{https://doi.org/10.1007/b13290}{Springer}, New York, 2009.

\bibitem{AG}
Andrews G.E., Garvan F.G., Ramanujan's ``lost'' notebook.~{VI}. {T}he mock
  theta conjectures,
  \href{https://doi.org/10.1016/0001-8708(89)90070-4}{\textit{Adv. Math.}}
  \textbf{73} (1989), 242--255.

\bibitem{AH}
Andrews G.E., Hickerson D., Ramanujan's ``lost'' notebook.~{VII}. {T}he sixth
  order mock theta functions,
  \href{https://doi.org/10.1016/0001-8708(91)90083-J}{\textit{Adv. Math.}}
  \textbf{89} (1991), 60--105.

\bibitem{Ba}
Bailey W.N., On the basic bilateral hypergeometric series~{${}_2\Psi_2$},
  \href{https://doi.org/10.1093/qmath/1.1.194}{\textit{Quart.~J. Math. Oxford
  Ser.}} \textbf{1} (1950), 194--198.

\bibitem{BKLMR}
Bajpai J., Kimport S., Liang J., Ma D., Ricci J., Bilateral series and
  {R}amanujan's radial limits,
  \href{https://doi.org/10.1090/S0002-9939-2014-12249-0}{\textit{Proc. Amer.
  Math. Soc.}} \textbf{143} (2015), 479--492,
  \href{http://arxiv.org/abs/2202.12141}{arXiv:2202.12141}.

\bibitem{B}
Berndt B.C., Ramanujan's notebooks. {P}art~{III},
  \href{https://doi.org/10.1007/978-1-4612-0965-2}{Springer}, New York, 1991.

\bibitem{BC}
Berndt B.C., Chan S.H., Sixth order mock theta functions,
  \href{https://doi.org/10.1016/j.aim.2007.06.004}{\textit{Adv. Math.}}
  \textbf{216} (2007), 771--786.

\bibitem{BFOR}
Bringmann K., Folsom A., Ono K., Rolen L., Harmonic {M}aass forms and mock
  modular forms: {T}heory and applications, \textit{Amer. Math. Soc. Colloq.
  Publ.}, Vol.~64, \href{https://doi.org/10.1090/coll/064}{American
  Mathematical Society}, Providence, RI, 2017.

\bibitem{BO1}
Bringmann K., Ono K., The~{$f(q)$} mock theta function conjecture and partition
  ranks, \href{https://doi.org/10.1007/s00222-005-0493-5}{\textit{Invent.
  Math.}} \textbf{165} (2006), 243--266.

\bibitem{BO2}
Bringmann K., Ono K., Dyson's ranks and {M}aass forms,
  \href{https://doi.org/10.4007/annals.2010.171.419}{\textit{Ann. of Math.}}
  \textbf{171} (2010), 419--449.

\bibitem{Ch}
Choi Y.-S., The basic bilateral hypergeometric series and the mock theta
  functions,
  \href{https://doi.org/10.1007/s11139-010-9269-7}{\textit{Ramanujan~J.}}
  \textbf{24} (2011), 345--386.

\bibitem{Dyson}
Dyson F.J., Some guesses in the theory of partitions, \textit{Eureka}
  \textbf{8} (1944), 10--15.

\bibitem{FG}
Frye J., Garvan F., Automatic proof of theta-function identities, in Elliptic
  {I}ntegrals, {E}lliptic {F}unctions and {M}odular {F}orms in {Q}uantum
  {F}ield {T}heory, \textit{Texts Monogr. Symbol. Comput.},
  \href{https://doi.org/10.1007/978-3-030-04480-0_10}{Springer}, Cham, 2019,
  195--258, \href{http://arxiv.org/abs/1807.08051}{arXiv:1807.08051}.

\bibitem{GR}
Gasper G., Rahman M., Basic hypergeometric series, 2nd ed., \textit{Encyclopedia Math.
  Appl.}, Vol.~96,
  \href{https://doi.org/10.1017/CBO9780511526251}{Cambridge University Pres}s,
  Cambridge, 2004.

\bibitem{GM2}
Gordon B., McIntosh R.J., Some eighth order mock theta functions,
  \href{https://doi.org/10.1112/S0024610700008735}{\textit{J.~London Math.
  Soc.}} \textbf{62} (2000), 321--335.

\bibitem{GM1}
Gordon B., McIntosh R.J., Modular transformations of {R}amanujan's fifth and
  seventh order mock theta functions, \href{https://doi.org/10.1023/A:1026299229509}{\textit{Ramanujan~J.}} \textbf{7} (2003), 193--222.

\bibitem{GM}
Gordon B., McIntosh R.J., A survey of classical mock theta functions, in
  Partitions, {$q$}-{S}eries, and {M}odular {F}orms, \textit{Dev. Math.},
  Vol.~23, \href{https://doi.org/10.1007/978-1-4614-0028-8_9}{Springer}, New
  York, 2012, 95--144.

\bibitem{GH}
Gu N.S.S., Hao L.-J., On some new mock theta functions,
  \href{https://doi.org/10.1017/s1446788718000368}{\textit{J.~Aust. Math.
  Soc.}} \textbf{107} (2019), 53--66.

\bibitem{H}
Hickerson D., A proof of the mock theta conjectures,
  \href{https://doi.org/10.1007/BF01394279}{\textit{Invent. Math.}} \textbf{94}
  (1988), 639--660.

\bibitem{HM}
Hickerson D., Mortenson E., Hecke-type double sums, {A}ppell--{L}erch sums, and
  mock theta functions,~{I},
  \href{https://doi.org/10.1112/plms/pdu007}{\textit{Proc. Lond. Math. Soc.}}
  \textbf{109} (2014), 382--422,
  \href{http://arxiv.org/abs/1208.1421}{arXiv:1208.1421}.

\bibitem{Hi}
Hikami K., Transformation formula of the ``second'' order {M}ock theta
  function, \href{https://doi.org/10.1007/s11005-005-0039-1}{\textit{Lett.
  Math. Phys.}} \textbf{75} (2006), 93--98,
  \href{http://arxiv.org/abs/math-ph/0604007}{arXiv:math-ph/0604007}.

\bibitem{HSZ}
Hu Q., Song H., Zhang Z., Third-order mock theta functions,
  \href{https://doi.org/10.1142/S1793042120500050}{\textit{Int.~J. Number
  Theory}} \textbf{16} (2020), 91--106.

\bibitem{HZ}
Hu Q., Zhang Z., On generalizations of theorems involving the third-order mock
  theta functions, \href{https://doi.org/10.1090/proc/14691}{\textit{Proc.
  Amer. Math. Soc.}} \textbf{148} (2020), 125--132.

\bibitem{K}
Kang S.-Y., Mock {J}acobi forms in basic hypergeometric series,
  \href{https://doi.org/10.1112/S0010437X09004060}{\textit{Compos. Math.}}
  \textbf{145} (2009), 553--565,
  \href{http://arxiv.org/abs/0806.1878}{arXiv:0806.1878}.

\bibitem{Mc}
Mc~Laughlin J., Mock theta function identities deriving from bilateral basic
  hypergeometric series, in Analytic {N}umber {T}heory, {M}odular {F}orms and
  {$q$}-{H}ypergeometric {S}eries, \textit{Springer Proc. Math. Stat.}, Vol.~221, \href{https://doi.org/10.1007/978-3-319-68376-8_29}{Springer}, Cham,
  2017, 503--531, \href{http://arxiv.org/abs/1906.11997}{arXiv:1906.11997}.

\bibitem{McI}
McIntosh R.J., Second order mock theta functions,
  \href{https://doi.org/10.4153/CMB-2007-028-9}{\textit{Canad. Math. Bull.}}
  \textbf{50} (2007), 284--290.

\bibitem{McI1}
McIntosh R.J., The~{$H$} and~{$K$} family of mock theta functions,
  \href{https://doi.org/10.4153/CJM-2011-066-0}{\textit{Canad.~J. Math.}}
  \textbf{64} (2012), 935--960.

\bibitem{Mo1}
Mortenson E., On the dual nature of partial theta functions and
  {A}ppell--{L}erch sums,
  \href{https://doi.org/10.1016/j.aim.2014.07.018}{\textit{Adv. Math.}}
  \textbf{264} (2014), 236--260,
  \href{http://arxiv.org/abs/1208.6316}{arXiv:1208.6316}.

\bibitem{Mo}
Mortenson E., Ramanujan's radial limits and mixed mock modular bilateral
  {$q$}-hypergeometric series,
  \href{https://doi.org/10.1017/S0013091515000425}{\textit{Proc. Edinb. Math.
  Soc.~(2)}} \textbf{59} (2016), 787--799,
  \href{http://arxiv.org/abs/1309.4162}{arXiv:1309.4162}.

\bibitem{R}
Ramanujan S., The lost notebook and other unpublished papers, Narosa Publishing
  House, New Delhi, 1988.

\bibitem{SW}
Song H., Wang C., Identities involving mock theta functions and theta
  functions, \href{https://doi.org/10.1090/proc/16343}{\textit{Proc. Amer.
  Math. Soc.}} \textbf{151} (2023), 3009--3022.

\bibitem{W1}
Watson G.N., The final problem: {A}n account of the mock theta functions,
  \href{https://doi.org/10.1112/jlms/s1-11.1.55}{\textit{J.~London Math. Soc.}}
  \textbf{11} (1936), 55--80.

\bibitem{W2}
Watson G.N., The mock theta functions~(2),
  \href{https://doi.org/10.1112/plms/s2-42.1.274}{\textit{Proc. London Math.
  Soc.}} \textbf{42} (1936), 274--304.

\bibitem{WYH}
Wei C., Yu Y., Hu Q., Two generalizations of {J}acobi's triple product identity
  and their applications,
  \href{https://doi.org/10.1007/s11139-022-00640-x}{\textit{Ramanujan~J.}}
  \textbf{60} (2023), 925--937.

\bibitem{ZS2}
Zhang W., Shi J., Terminating {$q$}-{W}hipple summation theorem with two
  integer parameters,
  \href{https://doi.org/10.1007/s11139-019-00208-2}{\textit{Ramanujan~J.}}
  \textbf{54} (2021), 137--146.

\bibitem{ZS}
Zhang Z., Song H., Some identities of certain basic hypergeometric series and
  their applications to mock theta functions,
  \href{https://doi.org/10.1007/s13226-022-00335-5}{\textit{Indian~J. Pure
  Appl. Math.}} \textbf{54} (2023), 1214--1225.

\bibitem{Z}
Zwegers S., Mock theta functions, Ph.D.~Thesis, {U}niversiteit Utrecht, 2002,
  \href{http://arxiv.org/abs/0807.4834}{arXiv:0807.4834}.

\bibitem{Z1}
Zwegers S., Mock {$\theta$}-functions and real analytic modular forms, in
  {$q$}-{S}eries with {A}pplications to {C}ombinatorics, {N}umber {T}heory, and
  {P}hysics ({U}rbana, {IL}, 2000), \textit{Contemp. Math.}, Vol.~291,
  \href{https://doi.org/10.1090/conm/291/04907}{American Mathematical Society},
  Providence, RI, 2001, 269--277.

\end{thebibliography}
